\subsection*{VI.22.P8: Factorization through a closed subscheme}
\begin{problem}\label{prob:vi22-p08}
Let $A\to B$ define $f:\Spec B\to\Spec A$ and let $Z=\Spec(A/I)$.  Prove $f$ factors through $Z$ iff $I$ maps to zero in $B$.
\end{problem}

\ifdefined\FullProblemDossiers
% VI22-FULL-P08
\paragraph{Why this problem matters.}
This is the universal property that turns a closed subscheme into the scheme-theoretic vanishing locus of equations.

\paragraph{Useful ideas before starting.}
Translate the factorization diagram contravariantly and invoke the universal property of the quotient ring.

\paragraph{Guided hints.}
\begin{enumerate}
\item Write the desired factorization $\Spec B\to\Spec(A/I)\to\Spec A$.
\item Reverse arrows to obtain $A\to A/I\to B$.
\item Apply the quotient universal property and note uniqueness.
\end{enumerate}

\begin{solution}
Let $\alpha:A\to B$ be the ring map corresponding to $f$.  A geometric
factorization
\[
 \Spec B\xrightarrow{\tilde f}\Spec(A/I)\longrightarrow\Spec A
\]
corresponds, by affine anti-equivalence, to a ring factorization
\[
 A\twoheadrightarrow A/I\xrightarrow{\tilde\alpha}B
\]
whose composite is $\alpha$.

By the universal property of the quotient, such a map $\tilde\alpha$ exists if and only if
$I\subseteq\ker\alpha$, equivalently every element of $I$ maps to zero in $B$.  When it exists,
it is unique because the quotient map $A\twoheadrightarrow A/I$ is surjective.  Hence the
scheme-theoretic factorization is also unique.

This is stronger than requiring only that the point-set image of $f$ lie in $V(I)$: pointwise
containment merely forces the image of $I$ to be nilpotent in $B$, not necessarily zero.
\end{solution}

\paragraph{What happened mathematically.}
Closed subschemes represent exact vanishing, while closed supports represent vanishing only modulo nilpotents.

\paragraph{Extensions.}
Construct a morphism whose point image lies in $V(I)$ but which does not factor through $\Spec(A/I)$.

\paragraph{Provenance.}
\textbf{Canonical VI/22 dossier developed from the AG5 T02 theory scope; no numbered legacy problem is claimed.}
\else
\paragraph{Provenance.} \textbf{Canonical VI/22 dossier developed from the AG5 T02 theory scope; no numbered legacy problem is claimed.}
\fi
